% Page 062 — page imprimée 59
% Georges Roth / Olivier Giran — fin du témoignage

Il se tourna vers la porte, et je lui dis encore pour terminer : « Je resterai avec vous jusqu’à la
dernière seconde... je vous accompagnerai comme un homme que je souhaiterai être ou avoir été
mon ami... ». Quelque chose d’intraduisible s’était passé entre nous….
Je dis encore une fois le « Notre Père » avec les trois condamnés. Ils prièrent tous trois à haute
voix. Une fois encore Olivier prononça d’une voix pleine de joie : « Car c’est à toi
qu’appartiennent, dans tous les siècles, le règne, la puissance et la gloire. Amen ».

Je l’entendis crier, au dernier moment dans le crépitement des salves « Vive la France ! ».
Moi-même, je le pleurai, lui que j’avais appris à respecter et à aimer en deux heures et je le pleure
encore maintenant en écrivant cette lettre.

Vous pouvez être fière d’un fils qui a su mourir fort et libre, vraiment fière. J’ai eu la certitude
qu’il était rempli de la pensée qu’il avait atteint dans la mort la plénitude de sa vie et qu’il ne
pourrait plus jamais la dépasser. Pour beaucoup on peut dire qu’ils ont dû mourir inachevés.
Olivier cependant est mort en sachant avec certitude que l’œuvre de sa vie était accomplie.

J’ai déjà prêché beaucoup depuis. Souvent, lorsque je parle de la certitude de la vie future, je
raconte à mes auditeurs la mort d’un homme appelé Giran qui traversa si superbement la mort
pour aller dans la vie éternelle. Et toujours, à nouveau, je vois que l’histoire de cette mort touche
le cœur de mes fidèles.

Vous voyez, sa mort elle-même touche beaucoup d’hommes qui ne le connaissaient pas et agit
comme un exemple qui les conduira à la vie éternelle.

\begin{flushright}
Hermann Hühn \ Pastor.
\end{flushright}

\vspace{1.5em}

\begin{tcolorbox}[colback=white,colframe=black]
\itshape
Après la mort de son mari et de son fils, Madame Giran, avec sa fille Claire, se résolut à
se séparer de la maison de Meyrueis.

Elle souhaitait vendre à des gens qui n’avaient pas été compromis dans la collaboration
ni dans le marché noir et c’est Gaston Poujol, leur plus proche voisin qui
reçut procuration pour s’occuper de cette vente.

M. Ernest Domergue, boucher du Vigan se porta acquéreur et revendit la maison
rapidement à M. Agrinier, boucher à Meyrueis.
\end{tcolorbox}
